% SPDX-FileCopyrightText: Copyright (c) 2016-2026 Objectionary.com
% SPDX-License-Identifier: MIT

\begin{proof}[Proof of \cref{thm:normalization-total}]
Consider any reduction sequence from a given expression under the rules of
  \cref{fig:reduction}; we show it is finite.

Away from \nameref{r:dot}, every rule strictly shrinks the term.
The \emph{collapsing} rules---\nameref{r:amiss}, \nameref{r:dc}, \nameref{r:dca},
  \nameref{r:dd}, \nameref{r:dl}, \nameref{r:miss}, \nameref{r:null},
  \nameref{r:over}, \nameref{r:overa}, and \nameref{r:stop}---rewrite their redex
  to the terminator $T$; \nameref{r:copy} and \nameref{r:stay} each discharge one
  application; and \nameref{r:alpha} turns a positional argument into a named one.
So only finitely many non-\nameref{r:dot} steps occur before either the term is
  exhausted or a \nameref{r:dot} redex is reached.

It remains to bound the \nameref{r:dot} steps.
\nameref{r:dot} replaces a dispatch \([[ B_1, \tau -> n, B_2 ]].\tau\) by the body
  \(n\) contextualized against \([[ B_1, B_2 ]]\) and decorated with the whole
  formation under \(\phiTerminal{\rho}\).
An infinite sequence would need infinitely many such steps, hence a
  \nameref{r:dot} step that reproduces a dispatch of an attribute on a formation
  that still binds it---a self-reproducing redex.
None can.
The body \(n\) is contextualized against the formation with the dispatched
  attribute removed, \([[ B_1, B_2 ]]\) (\cref{sec:contextualization}), so every
  \(\phiTerminal{\xi}\) in \(n\) that denotes the enclosing formation becomes a
  formation binding neither \(\tau\) nor \(\phiTerminal{\rho}\); a self-reference
  that reaches for \(\tau\) again, whether directly as \(\phiTerminal{\xi}.\tau\)
  or through the parent as \(\phiTerminal{\xi}.\phiTerminal{\rho}.\tau\),
  dispatches a missing attribute and collapses to $T$ by \nameref{r:stop},
  \nameref{r:null}, \nameref{r:miss}, or \nameref{r:dd}.
The copy of the formation stored under \(\phiTerminal{\rho}\) is reachable only by
  dispatching \(\phiTerminal{\rho}\), again through such a chain, which dies the
  same way.

The calculus expresses recursion in only two further ways, and neither is a
  reduction.
Firing the atom that guards a recursive object---the base-case test that decides
  whether to recur---belongs to evaluation (\cref{sec:evaluation}): a formation
  carrying an unfired $L$-asset matches no rule of \cref{fig:reduction}
  (\nameref{r:stop} excludes it), so the recursive branch never unfolds under
  normalization.
Following the \(\phiTerminal{\phi}\)-decoration to an inherited attribute belongs
  to morphing (\cref{sec:morphing}): no reduction descends into
  \(\phiTerminal{\phi}\), so dispatching an absent attribute of a decorated
  formation is irreducible rather than chasing a possibly cyclic chain.

Every dispatch a term exposes is thus drawn from its finite structure and consumed
  once, so \nameref{r:dot} fires finitely often and the sequence is finite.
Every expression therefore has a normal form, unique by
  \cref{thm:normalization-confluent}.
This argument is not yet mechanized; extending the Lean~4 development that
  certifies \cref{thm:normalization-confluent} to strong normalization remains
  future work.
\end{proof}
